\chapter{Schéma zapojení a popis funkce}
\label{chap:napajeni}
\section{Společné napájení a výstup k~servu}
Servo i řízení napájíme stabilizovaným zdrojem 5 V / alespoň 4 A. Za pomalou pojistkou F1 T 3,15 A se větev rozdělí na samostatný silový přívod serva a přívod řízení. Proud motoru nesmí procházet Arduinem. Zemní návraty se setkají v~G0 u~přívodu zdroje; společná zem je referencí signálu.

Úplné zapojení s~Arduinem je na obr.~\ref{fig:knob-priklad}. Analogová varianta používá stejný výstupní blok, pouze místo D9 připojíme OUT časovače. Oba výstupy nikdy nespojujeme. R6 = 330~$\Omega$ omezuje proud v~signálové cestě a R7 = 100 k$\Omega$ drží nepřipojený vstup u~nuly. Nejde o~galvanické oddělení ani úplnou ochranu proti přepětí.

C9 = 1000 µF a C10 = 100 nF umístíme u~konektoru serva. První pomáhá pokrýt krátké změny odběru, druhý rychlé složky rušení. Kondenzátor ale nenahradí zdroj: při proudu 1 A po dobu 1 ms by napětí na samotných 1000 µF kleslo o~$\Delta U=I\Delta t/C=1$ V. Při zablokování trvá odběr dlouhodobě a je nutné pohon zastavit nebo odpojit.

Na servu musí při rozběhu zůstat alespoň 4,8 V. Z~jmenovitých 5 V tak zbývá jen 0,2 V pro toleranci zdroje, pojistku, kontakty a kabely. Sériovou ochrannou diodu kvůli úbytku nepoužíváme; polaritu zajišťuje správně zapojený nezáměnný konektor. Elektronickou ochranu proti přepólování návrh nemá. F1 chrání přívod, ne trvale zablokovaný motor. Výběr musí odpovídat vodičům a zkratovému proudu zdroje.

Na desce vedeme rychlé návratové proudy krátkou cestou a signál nad souvislou referenční plochou. Napájecí vodiče serva držíme blízko sebe, blokování u~konektoru a citlivý vstup potenciometru mimo silové vedení. Nízká frekvence povelu 50 Hz neodstraňuje vliv rychlých hran. Tyto zásady navazují na doporučení TI a Analog Devices.~\cite{tiemi,mt031}

Přívody zajistíme proti tahu a odpojovač napájení umístíme do dosahu obsluhy.

\section{Zapojení s~jedním NE555}
Schéma ukazuje všech osm vývodů a tři kontakty potenciometru P1. Jezdec W je spojen s~koncem B; při výpadku kontaktu tak zůstává zapojen celý odpor dráhy. Pulz se ovšem prodlouží, takže toto spojení nenahrazuje omezení pohybu. R1 omezuje proud i při P1 nastaveném na nulu. D1 má anodu na vývodu 7 a katodu na společném uzlu vývodů 2/6.

\noindent\begin{minipage}{\textwidth}\centering
\begin{tikzpicture}[transform shape]
	% Paths, nodes and wires:
	\draw[line width=1.2pt,fill=light] (7, 2.5) -- (11.25, 2.5) -| (11.25, 7) -- (7, 7) -| (7, 2.5);
	\node[shape=rectangle, minimum width=0.965cm, minimum height=0.365cm] at (9.75, 5.05){} node[anchor=center, align=center, text width=0.929cm, inner sep=1pt] at (9.75, 5.05){\small U1};
	\node[shape=rectangle, minimum width=1.765cm, minimum height=0.465cm] at (9.85, 4.5){} node[anchor=center, align=center, text width=1.729cm, inner sep=1pt] at (9.85, 4.5){\small NE555P};
	\draw (7, 6.3) -- (3, 6.3);
	\draw (7, 4.5) -- (6.2, 4.5) -| (6.2, 3.4);
	\draw (7, 3.4) -- (3, 3.4);
	\node[shape=rectangle, minimum width=1.465cm, minimum height=0.315cm] at (7.9, 6.3){} node[anchor=west, align=left, text width=1.429cm, inner sep=1pt] at (7.15, 6.3){\scriptsize 7 DISCH};
	\node[shape=rectangle, minimum width=1.465cm, minimum height=0.315cm] at (7.9, 4.5){} node[anchor=west, align=left, text width=1.429cm, inner sep=1pt] at (7.15, 4.5){\scriptsize 6 THRES};
	\node[shape=rectangle, minimum width=1.465cm, minimum height=0.315cm] at (7.9, 3.4){} node[anchor=west, align=left, text width=1.429cm, inner sep=1pt] at (7.15, 3.4){\scriptsize 2 TRIG};
	\draw (11.25, 4.75) |- (12.3, 4.8);
	\node[shape=rectangle, minimum width=2.465cm, minimum height=0.815cm] at (13.7, 4.8){} node[anchor=west, align=left, text width=2.429cm, inner sep=1pt] at (12.45, 4.8){\small 3 OUT\\signál pro servo};
	\draw (8, 7) |- (7.75, 11) -| (9.5, 7.5) -| (9.5, 7);
	\node[shape=rectangle, minimum width=0.965cm, minimum height=0.265cm] at (8, 6.75){} node[anchor=center, align=center, text width=0.929cm, inner sep=1pt] at (8, 6.75){\scriptsize 8 VCC};
	\node[shape=rectangle, minimum width=1.865cm, minimum height=0.115cm] at (9.8, 6.675){} node[anchor=center, align=center, text width=1.829cm, inner sep=1pt] at (9.8, 6.675){\scriptsize 4 RESET};
	\draw (8, 2.5) -| (8, 1.5);
	\node[ground] at (8, 0.5){};
	\node[shape=rectangle, minimum width=1.165cm, minimum height=0.315cm] at (8, 2.75){} node[anchor=center, align=center, text width=1.129cm, inner sep=1pt] at (8, 2.75){\scriptsize 1 GND};
	\draw (9.5, 2.5) to[capacitor] (9.5, 0.6);
	\node[ground] at (9.5, 0.5){};
	\node[shape=rectangle, minimum width=1.265cm, minimum height=0.315cm] at (9.5, 2.75){} node[anchor=center, align=center, text width=1.229cm, inner sep=1pt] at (9.5, 2.75){\scriptsize 5 CONT};
	\node[shape=rectangle, minimum width=1.965cm, minimum height=0.465cm] at (11, 1.55){} node[anchor=center, align=center, text width=1.929cm, inner sep=1pt] at (11, 1.55){\small C7 10 nF};
	\node[shape=rectangle, minimum width=1.465cm, minimum height=0.465cm] at (3, 11.789){} node[anchor=center, align=center, text width=1.429cm, inner sep=1pt] at (3, 11.789){\small +5 V};
	\draw (3, 11) to[european resistor] (3, 9.2);
	\draw (3, 9.2) -| (3, 8.8);
	\node[shape=rectangle, minimum width=2.465cm, minimum height=0.465cm] at (1.35, 10.1){} node[anchor=center, align=center, text width=2.429cm, inner sep=1pt] at (1.35, 10.1){\small R1 3,3 kΩ};
	\draw (2.75, 8.8) -- (3.25, 8.8) -| (3.25, 7.3) -- (2.75, 7.3) -| (2.75, 8.8);
	\node[shape=rectangle, minimum width=2.265cm, minimum height=0.815cm] at (1.3, 8.05){} node[anchor=center, align=center, text width=2.229cm, inner sep=1pt] at (1.3, 8.05){\small P1\\25 kΩ lin.};
	\node[shape=rectangle, minimum width=0.365cm, minimum height=0.265cm] at (2.65, 9){} node[anchor=center, align=center, text width=0.329cm, inner sep=1pt] at (2.65, 9){\scriptsize A};
	\node[shape=rectangle, minimum width=0.365cm, minimum height=0.265cm] at (2.65, 7.05){} node[anchor=center, align=center, text width=0.329cm, inner sep=1pt] at (2.65, 7.05){\scriptsize B};
	\draw[-latex] (4.5, 8.05) -- (3.25, 8.05);
	\node[shape=rectangle, minimum width=0.415cm, minimum height=0.315cm] at (4.85, 8.05){} node[anchor=center, align=center, text width=0.379cm, inner sep=1pt] at (4.85, 8.05){\scriptsize W};
	\draw (4.5, 8.05) -| (4.5, 7) -- (3, 7);
	\draw (3, 7.3) -| (3, 6.3);
	\draw (3, 6.3) to[european resistor] (3, 3.4);
	\node[shape=rectangle, minimum width=2.565cm, minimum height=0.465cm] at (1.3, 4.85){} node[anchor=center, align=center, text width=2.529cm, inner sep=1pt] at (1.3, 4.85){\small R2 221 kΩ};
	\draw (5.1, 6.3) to[empty diode] (5.1, 3.4);
	\node[shape=rectangle, minimum width=1.165cm, minimum height=0.615cm] at (5.9, 5.551){} node[anchor=center, align=center, text width=1.129cm, inner sep=1pt] at (5.9, 5.551){\scriptsize D1\\1N4148};
	\node[shape=rectangle, minimum width=0.365cm, minimum height=0.265cm] at (4.7, 5.8){} node[anchor=center, align=center, text width=0.329cm, inner sep=1pt] at (4.7, 5.8){\scriptsize A};
	\node[shape=rectangle, minimum width=0.365cm, minimum height=0.265cm] at (4.7, 4){} node[anchor=center, align=center, text width=0.329cm, inner sep=1pt] at (4.7, 4){\scriptsize K};
	\draw (3, 3.4) to[capacitor] (3, 1.2);
	\node[ground] at (3, 0.5){};
	\node[shape=rectangle, minimum width=2.565cm, minimum height=0.465cm] at (1.3, 2.3){} node[anchor=center, align=center, text width=2.529cm, inner sep=1pt] at (1.3, 2.3){\small C1 120 nF};
	\node[circ] at (3, 7){};
	\node[circ] at (3, 6.3){};
	\node[circ] at (5.1, 6.3){};
	\node[circ] at (3, 3.4){};
	\node[circ] at (5.1, 3.4){};
	\node[circ] at (6.2, 3.4){};
	\draw (10.75, 10.95) to[capacitor] (10.75, 8.75);
	\node[ground] at (10.75, 8.75){};
	\draw (13.45, 10.95) to[capacitor] (13.45, 8.75);
	\node[ground] at (13.45, 8.75){};
	\node[shape=rectangle, minimum width=2.165cm, minimum height=0.465cm] at (11.65, 9.25){} node[anchor=center, align=center, text width=2.129cm, inner sep=1pt] at (11.65, 9.25){\small C5 100 nF};
	\node[shape=rectangle, minimum width=2.165cm, minimum height=0.465cm] at (14.4, 9.25){} node[anchor=center, align=center, text width=2.129cm, inner sep=1pt] at (14.4, 9.25){\small C6 10 µF};
	\node[shape=rectangle, minimum width=0.365cm, minimum height=0.365cm] at (13, 10.2){} node[anchor=center, align=center, text width=0.329cm, inner sep=1pt] at (13, 10.2){\small +};
	\node[shape=rectangle, minimum width=4.365cm, minimum height=0.765cm] at (12.056, 7.519){} node[anchor=center, align=center, text width=4.329cm, inner sep=1pt] at (12.056, 7.519){\scriptsize C5 umístit co nejblíže\\napájecím vývodům NE555};
	\draw (8, 1.5) -| (8, 0.5);
	\draw (9.5, 0.6) -| (9.5, 0.5);
	\draw (3, 1.2) -- (3, 0.5);
	\draw (3, 11) -| (13.45, 10.95);
	\draw (10.75, 10.95) -| (10.75, 11);
	\node[vcc] at (3, 11){};
	\node[circ] at (8, 11){};
	\node[circ] at (9.5, 11){};
	\node[circ] at (10.75, 11){};
	\node[circ] at (3, 11){};
\end{tikzpicture}
\captionsetup{hypcap=false}\captionof{figure}{Generátor povelu s~jedním NE555.}
\label{fig:ne555}
\end{minipage}\medskip

C5 = 100 nF blokuje napájení u~vývodů 8 a 1, C6 = 10 µF doplňuje místní zásobu energie a C7 = 10 nF filtruje CONT. RESET je připojen na +5 V. Výstup 3 vede přes R6 na signál serva, zem časovače samostatně do G0. NE555 vytváří pouze povel, nenapájí motor.~\cite{ne555}

\section{Společná zem a návrat proudu}
Skutečný zemní vodič má odpor i indukčnost. Proud motoru ve společném úseku vytváří rušivé napětí
\begin{equation}
 u_\mathrm{GND}=R_\mathrm{spol}i+L_\mathrm{spol}\frac{\mathrm{d}i}{\mathrm{d}t}.
\end{equation}
Například 2 A přes 0,1 $\Omega$ způsobí úbytek 0,2 V. Pokud se přičte k~měření potenciometru, naruší povel. Proto zem P2 a C11 připojíme u~Arduina a odtud vedeme samostatnou větev do G0. Silový návrat serva vede přímo ke zdroji.

\noindent\begin{minipage}{\textwidth}\centering
\resizebox{.90\textwidth}{!}{\begin{circuitikz}[x=1cm,y=1cm]
\node[anchor=west,font=\bfseries\small] at (-1.2,2) {a) Proud serva protéká společným úsekem};
\node[blok,text width=20mm] (s) at (0,1) {Zdroj\\záporný pól};
\node[blok,text width=23mm] (a) at (5,1) {Arduino\\GND};
\node[blok,text width=20mm] (m) at (11,1) {Servo\\GND};
\draw (s.south)--(0,-.3);
\draw[very thick,red!60!black] (0,-.3) to[R,l=$R_\mathrm{spol}$] (5,-.3);
\draw[very thick,accent] (5,-.3)--(11,-.3)--(m.south);
\draw (a.south)--(5,-.3);\fill (5,-.3) circle(2pt);
\draw[tok,accent] (10,-.85)--node[below]{$I_\mathrm{servo}$}(6,-.85);
\draw[tok,red!60!black] (4.5,-.85)--node[below]{$\Delta u=R_\mathrm{spol}I_\mathrm{servo}$}(.5,-.85);
\node[anchor=west,font=\footnotesize] at (-1.2,-1.7) {Úbytek na společném vodiči posouvá zem Arduina vůči zdroji.};
\begin{scope}[yshift=-5.1cm]
\node[anchor=west,font=\bfseries\small] at (-1.2,2) {b) Samostatný návrat serva do společného bodu};
\node[blok,text width=20mm] (s2) at (0,1) {Zdroj\\záporný pól};
\node[blok,text width=23mm] (a2) at (5,1) {Arduino\\GND};
\node[blok,text width=20mm] (m2) at (11,1) {Servo\\GND};
\draw (s2.south)--(0,-1);
\draw (a2.south)--(5,-.1)--(0,-.1);
\draw[very thick,accent] (m2.south)--(11,-1)--(0,-1);
\fill (0,-.1) circle(2pt);\node[left] at (0,-.1) {G0};
\draw[tok,accent] (9,-1.4)--node[below]{$I_\mathrm{servo}$ mimo zemní větev Arduina}(2,-1.4);
\node[above,font=\footnotesize] at (3,-.1) {návrat řídicí části};
\end{scope}
\end{circuitikz}
}
\captionsetup{hypcap=false}\captionof{figure}{Společný zemní úsek a oddělené návraty.}
\label{fig:zeme}
\end{minipage}\medskip

Do úbytku napětí započítáme oba přívody, pojistku, konektory i kabel serva. Napětí změříme pod zatížením přímo na konektoru.

\section{Zemní smyčka a rušení}
Zemní smyčka vzniká, když další vodivá cesta uzavře okruh. Typickým příkladem je spojení počítače, Arduina a laboratorního zdroje přes USB, GND a ochrannou zem PE. Rozdíl potenciálů nebo indukce pak může vyvolat vyrovnávací proud. Společná impedance naproti tomu ruší i bez druhé cesty; nejde o~stejný jev.~\cite{osak}

\noindent\begin{minipage}{\textwidth}\centering
\resizebox{.90\textwidth}{!}{\begin{tikzpicture}[x=1cm,y=1cm,font=\small]
\node[blok,text width=29mm] (pc) at (0,3.4) {Stolní počítač\\USB GND};
\node[blok,text width=29mm] (uno) at (6.5,3.4) {Arduino\\GND};
\node[blok,text width=29mm] (src) at (6.5,0) {Laboratorní zdroj\\záporný pól};
\node[blok,text width=29mm] (pe) at (0,0) {Společný PE\\ochranný vodič};
\draw[thick] (pc.east)--node[above]{zem kabelu USB}(uno.west);
\draw[thick] (uno.south)--node[right]{přes G0}(src.north);
\draw[thick,dashed,accent] (src.south)--(6.5,-1.1)--node[below,align=center]{spojení záporného pólu s PE\\pouze pokud je ve zdroji přítomné}(0,-1.1)--(pe.south);
\draw[thick,dashed,accent] (pe.north)--node[left,align=right,font=\footnotesize]{vazba USB GND\\na kostru / PE}(pc.south);
\draw[tok,red!60!black] (1.9,2.6)--(4.8,2.6)--(4.8,.8)--(1.9,.8)--(1.9,2.2);
\node[align=center,font=\footnotesize] at (3.35,1.7) {možný proud\\zemní smyčkou};
\node[anchor=north west,text width=35mm,align=left,font=\footnotesize] at (8.5,3) {Plné čáry: zemní propojení zařízení.\\[3mm]Čárkovaně: další vazby na PE, které je nutné ověřit.};
\end{tikzpicture}
}
\captionsetup{hypcap=false}\captionof{figure}{Možná zemní smyčka přes USB a PE.}
\end{minipage}\medskip

Nejprve ověříme skutečná propojení, včetně vazby záporného pólu zdroje na PE a připojených měřicích přístrojů. Pro samostatný provoz USB odpojíme. Pokud je komunikace nutná a ruší, lze navrhnout izolované datové rozhraní. \textbf{Ochranný vodič PE se kvůli rušení nikdy neodpojuje.} Galvanické oddělení by vyžadovalo oddělené napájecí domény i přenos povelu přes izolátor.

